\chapter{التقارير والذكاء الاصطناعي}
\label{ch:reports}

تتعلم كيف تصل إلى أي نتيجة أو تقرير تحتاجه من مكان واحد، وأين يساعدك الذكاء الاصطناعي في يومك وما حدوده. ثم تختبر نفسك في تقييم قصير يغطي الحقيبة كلها.

\section{أي نتيجة تحتاجها، من مكان واحد}
\label{sec:rep-ways}
\begin{rcslide}{أي نتيجة تحتاجها، من مكان واحد}
\sphase{التقارير}
\slead{لأن عمل المؤسسة كله في كيان، تصل إلى الجواب بأربع طرق، من الأسرع إلى الأعمق.}
\scard{tasks}{طرق العرض والمرشحات}{تصفّي قائمة لترى ما تحتاجه الآن.}
\scard{chart}{لوحات المعلومات}{أرقام دائمة يتابعها الفريق.}
\scard{ai}{تقارير الذكاء الاصطناعي}{تسأل بلغتك، فتحصل على تقرير مكتوب.}
\scard{flow}{لوحات الذكاء الاصطناعي}{تصف اللوحة، فتُبنى لك.}
\stake{لا تجمع الأرقام من برامج متفرقة؛ اختر الطريقة التي تناسب سؤالك.}
\end{rcslide}

\section{طرق العرض والمرشحات}
\label{sec:rep-views}
\begin{rcslide}{طرق العرض والمرشحات}
\sphase{التقارير}
\slead{أسرع تقرير هو قائمة مصفّاة: لا تبني شيئاً، بل تغيّر طريقة رؤية ما هو موجود.}
\scard{tasks}{قائمة}{لمراجعة التفاصيل سجلاً سجلاً.}
\scard{projects}{لوحة}{لتوزيع العمل حسب الحالة.}
\scard{paper}{طرق عرض محفوظة}{تعود إليها بضغطة، وقد تظهر في قائمة التطبيق مع عدد سجلاتها.}
\stake{طريقة العرض لا تغيّر البيانات نفسها؛ تغيّر ما تراه منها فقط.}
\slink{main}{طرق العرض في المهام}{ar/tasks-views.html}
\end{rcslide}

\section{لوحات المعلومات}
\label{sec:rep-dash}
\begin{rcslide}{لوحات المعلومات}
\sphase{التقارير}
\slead{لوحة تُبنى مرة، ويتابعها الفريق كل يوم بأرقام من السجلات نفسها.}
\scard{chart}{المرشحات}{تغيّر الفترة أو الفرع، فتتحدث الأرقام.}
\scard{sheet}{تصدير الجدول}{تأخذ بيانات جدول من اللوحة.}
\scard{paper}{تقرير PDF}{تصدّر اللوحة كلها تقريراً.}
\scard{lock}{بحسب صلاحياتك}{يرى كل مشاهد ما تسمح به صلاحياته.}
\stake{التقرير الشهري الذي كان يُجمع يدوياً صار لوحة مفتوحة طوال الشهر.}
\slink{builder}{عرض اللوحات وتصديرها}{ar/builder-dashboards-viewing.html}
\end{rcslide}

\section{اسأل بلغتك}
\label{sec:rep-ask}
\begin{rcslide}{اسأل بلغتك}
\sphase{الذكاء الاصطناعي}
\slead{في «تقارير الذكاء الاصطناعي» تكتب سؤالك كما تقوله لزميل، فتحصل على تقرير مكتوب وأرقام. وفي «لوحات معلومات الذكاء الاصطناعي» تصف لوحة فتُبنى لك.}
\squote{كم زيارة أُلغيت هذا الشهر، ولماذا؟}
\scard{lock}{مساحة خاصة بك}{تستخدم فقط البيانات التي تستطيع قراءتها.}
\scard{people}{المشاركة}{لتشارك لوحة مع الفريق، انشرها من صفحة «الصفحات».}
\stake{السؤال الذي كان يأخذ يومين من الجمع صار يأخذ دقائق.}
\slink{builder}{تقارير الذكاء الاصطناعي}{ar/builder-ai-reports.html}
\slink{builder}{لوحات معلومات الذكاء الاصطناعي}{ar/builder-ai-dashboards.html}
\end{rcslide}

\section{كيف تسأل سؤالاً جيداً؟}
\label{sec:rep-good-q}
\begin{rcslide}{كيف تسأل سؤالاً جيداً؟}
\sphase{الذكاء الاصطناعي}
\slead{كلما كان سؤالك محدداً، كان الجواب أدق وأسرع.}
\scompare{سؤال عام}{«كيف حال العمل؟» \\ «أعطني تقريراً عن الطلبات»}{سؤال محدد}{«كم طلب صيانة فُتح في تشرين الثاني، وكم منها متأخر؟» \\ «قارن عدد الكشوف بين فرعي دمشق وحلب هذا الشهر»}
\scard{flow}{حدّد الفترة}{هذا الشهر؟ هذا الربع؟}
\scard{building}{حدّد النطاق}{أي فرع؟ أي وحدة؟ أي نوع سجلات؟}
\scard{chart}{حدّد المقياس}{عدد؟ نسبة؟ متوسط مدة؟}
\stake{راجع الجواب قبل أن تبني عليه قراراً، كما تراجع أي تقرير يصلك.}
\end{rcslide}

\section{الذكاء الاصطناعي في يومك}
\label{sec:rep-ai-day}
\begin{rcslide}{الذكاء الاصطناعي في يومك}
\sphase{الذكاء الاصطناعي}
\slead{لا يقتصر الذكاء الاصطناعي في كيان على التقارير؛ تجده حيث تعمل.}
\scard{mail}{في البريد}{ملخص الرسالة، وكشف الاحتيال والرسائل المزعجة.}
\scard{chat}{في المحادثات}{تفريغ الرسائل الصوتية نصاً، والبوتات.}
\scard{form}{في الكتابة}{مساعد الكتابة، وملخص الإشعارات.}
\scard{entity}{في البناء}{يقترح حقول الكيان من وصفك، وأنت تراجع قبل الحفظ.}
\scard{tasks}{في الرقابة}{نتائج الذكاء الاصطناعي: ملاحظات تحتاج إلى قرار.}
\stake{في كل هذه الأماكن يقترح الذكاء الاصطناعي، والقرار لك.}
\end{rcslide}

\section{الحدود والأمان}
\label{sec:rep-limits}
\begin{rcslide}{الحدود والأمان}
\sphase{الذكاء الاصطناعي}
\slead{يضبط مسؤول المؤسسة أين يعمل الذكاء الاصطناعي، ولمن، وبأي ميزانية.}
\scard{lock}{ضمن صلاحياتك}{لا يقرأ إلا ما يحق لك رؤيته.}
\scard{building}{توفر الذكاء الاصطناعي}{كل ميزة تُفعَّل أو تُوقف للمؤسسة أو لأعضاء محددين.}
\scard{chart}{الحدود}{لكل عضو وللمؤسسة رصيد؛ وإذا نفد توقفت الطلبات الجديدة.}
\stake{إذا لم تجد ميزة ذكاء اصطناعي يستخدمها زميلك، فقد تكون موقوفة لك، أو نفد رصيدك.}
\slink{main}{توفر الذكاء الاصطناعي}{ar/settings-ai-availability.html}
\slink{main}{حدود الذكاء الاصطناعي}{ar/settings-ai-limits.html}
\end{rcslide}

\section{التقييم الختامي: الفكرة والأدوات الجاهزة}
\label{sec:rep-quiz-1}
\begin{rcslide}{التقييم الختامي: الفكرة والأدوات الجاهزة}
\sphase{التقييم}
\slead{أجب في ذهنك أو على ورقة، ثم افتح الإجابة.}
\squiz{ما الفائدة الكبرى من جمع عمل المؤسسة في كيان؟}{برنامج أجمل من البرامج الحالية \\ عمل المؤسسة كله في مكان واحد ومترابط، فتصل إلى أي نتيجة من هنا \\ الاستغناء عن الموظفين}{2}{الترابط هو ما يجعل أي تقرير ممكناً من مكان واحد، بدل الجمع من برامج متفرقة.}
\squiz{طلب زميلك منك عملاً في المحادثات. ما الأفضل؟}{تتذكره وتنفذه حين تتفرغ \\ تنسخه إلى ملف Excel \\ «إنشاء مهمة» من الرسالة، فيصبح له مسؤول وموعد}{3}{المهمة تحفظ العمل ومسؤوله وموعده وحالته، فلا يضيع بين الرسائل.}
\squiz{بأي صفة تعمل في المراسلات؟}{باسمك الشخصي \\ باسم المنصب الذي تشغله \\ باسم وحدتك فقط}{2}{الصفة هي المنصب الذي تعمل باسمه، وتغييرها لا يمنحك صلاحية جديدة.}
\stake{إن أخطأت في سؤال، عد إلى الوحدة التي يشرحها.}
\end{rcslide}

\section{التقييم الختامي: البناء والتحول والتقارير}
\label{sec:rep-quiz-2}
\begin{rcslide}{التقييم الختامي: البناء والتحول والتقارير}
\sphase{التقييم}
\slead{ثلاثة أسئلة أخيرة عن أدوات البناء ونقل العمل والتقارير.}
\squiz{بأي أداة تبدأ حين تنقل دفتراً ورقياً إلى كيان؟}{لوحة المعلومات \\ الكيان \\ الإجراء الداخلي}{2}{الكيان هو الأساس: النموذج يكتب فيه، والإجراء يبدأ منه، واللوحة تحسب منه.}
\squiz{عدّلت نموذجاً في أدوات البناء ولم يرَ زملاؤك التعديل. لماذا؟}{التعديل ما زال مسودة لم تُنشر \\ النموذج معطّل \\ زملاؤك لا يملكون حساباً}{1}{كل تعديل يبقى مسودة لا يراها غيرك حتى تراجعه وتنشره.}
\squiz{سألت الذكاء الاصطناعي عن بيانات لا تملك صلاحية رؤيتها. ماذا يحدث؟}{يعرضها لك كاملة \\ لا يستخدم إلا البيانات التي تستطيع قراءتها \\ يرسل طلباً إلى المدير}{2}{تعمل التحليلات بالذكاء الاصطناعي ضمن صلاحياتك، فلا تكشف ما لا يحق لك رؤيته.}
\stake{أنهيت الحقيبة. الخطوة التالية: اختر تطبيقاً عملياً وطبّقه خطوة خطوة.}
\slink{inspection}{تطبيق عملي: الكشف على المنشآت}{ar/index.html}
\slink{correspondence}{تطبيق عملي: المراسلات}{ar/index.html}
\slink{sop}{تطبيق عملي: الإجراءات الداخلية}{ar/index.html}
\end{rcslide}
